\section{Discussion}

\subsection{What is difficult to copy}

The surface architecture of \systemname is intentionally modular. A capable
synthesizer, several producers, and a router can be reproduced. This is not the
primary moat. The compounding asset is the joint empirical model required to
make the architecture effective:
\begin{itemize}
  \item a versioned capability map over model, role, task, and state;
  \item observed redundancy and error correlation among evidence sources;
  \item a representation of unresolved claims, readiness, and risk;
  \item an allocation and stopping policy trained on task outcomes; and
  \item a dataset connecting orchestration decisions to long-horizon success.
\end{itemize}

More usage can improve each component. Better allocation produces better data
about where models are useful; better data produces better allocation. The
feedback loop is defensible only when outcome attribution, privacy, evaluation,
and rollback are first-class system properties.

\subsection{Why singular authority matters}

Many multi-agent systems treat autonomy as a source of collective capability.
For a model exposed as one API identity, autonomy also creates a consistency
problem. Agents can act from incompatible assumptions, issue duplicate tool
calls, or complete against different interpretations of success. \systemname
separates epistemic plurality from operational plurality: many processes may
search, but one process owns the next externally visible commitment.

This boundary is particularly valuable when the system is embedded inside
another agent. From the caller's perspective, \systemname behaves as one model
with one tool protocol and one conversation state. Internally, it can purchase
as much plural evidence as the decision warrants without exposing a swarm's
coordination burden to the client.

\subsection{Limits of the information-theoretic framing}

The information-acquisition view is a guide, not a claim that natural-language
tasks admit a clean finite hypothesis space. Semantic clusters can be wrong.
Model self-confidence is poorly calibrated. Creative and preference-sensitive
tasks may have no unique latent truth. Evidence sources share data, training
methods, and failure modes that are difficult to observe. The value of a call
can become clear only many steps later.

Consequently, a universal scalar ``uncertainty score'' would be misleading.
Solve, Explore, Variant, and Critique require different state representations
and success signals. Initial policies should remain typed, auditable, and
conservative where actions are consequential. Learned policies should report
calibration and out-of-distribution behavior rather than only average reward.

\subsection{A research agenda for the new scaling dimension}

Three questions determine whether allocation intelligence becomes a durable
scaling dimension.

First, can conditional contribution be predicted well enough to outperform
fixed test-time scaling across changing model pools? Second, can a system learn
to stop without losing rare but consequential corrections? Third, can
trajectory outcomes provide sufficiently local credit to improve allocation
without overfitting to a benchmark or existing policy?

Answering these questions requires scaling experiments, not only endpoint
comparisons. Future evaluations should vary worker capability, pool diversity,
call budget, workflow length, and task distribution. A genuine scaling result
would show that allocation improvements shift the system frontier repeatedly,
not merely at one hand-tuned operating point.

